\documentclass[12pt]{article}
\usepackage[utf8]{inputenc}
\usepackage{amsmath, amssymb}
\usepackage{geometry}
\geometry{margin=1in}

\title{On the Inherent Limits of Classical Mathematical Frameworks for Structural Problems: An Entropy-Based Perspective}
\author{Masaomi Chiba}
\date{2025-12-17}

\begin{document}
\maketitle

\section*{Abstract}
We present a formal argument highlighting the intrinsic limitations of existing mathematical frameworks, particularly classical and contemporary axiom systems, in addressing structural problems such as the $\mathrm{ABC}$ conjecture and the Collatz conjecture. By framing these systems as structural fixed points of minimal reasoning entropy, we demonstrate that the true proofs of such conjectures require operations of ultimate logical refinement, which cannot be captured by any finite combination of conventional mathematical vocabulary.

\section{Introduction}
All extant mathematical theories, defined by their axiom sets and formal rules, can be viewed as structural fixed points wherein the reasoning entropy is locally minimized. Within this view, any proof conducted in these frameworks is inherently limited by the operational vocabulary and permissible transformations encoded within the theory. \textbf{Here, reasoning entropy is understood as the physical information cost associated with logical state transitions, consistent with Landauer's principle.}

\section{Structural Problems Beyond Classical Fixed Points}
Consider structural problems $P$ such as the $\mathrm{ABC}$ conjecture or the Collatz conjecture. These problems possess inherent contradictions between additive and multiplicative or recursive structures, which cannot be fully resolved within the confines of a fixed point of minimal reasoning entropy. 

Specifically, the $\mathrm{ABC}$ conjecture epitomizes a fundamental \textbf{Non-Middle Fallacy (NMF)} within classical arithmetic: the forced symmetry between additive and multiplicative structures. In a static axiom system, the equality in $a + b = c$ is treated as a rigid constraint, yet it fails to account for the divergent reasoning entropy required to reconcile it with the prime-factor (multiplicative) density of $abc$. Formally, let $F$ denote the structural fixed point associated with a given axiom system $\mathcal{A}$. Then, for any proof $\pi$ within $\mathcal{A}$, the reachability condition:
\[
\pi: F \rightarrow P
\]
cannot be satisfied if $P$ requires a dynamic transformation between these universes that lies beyond the static vocabulary of $\mathcal{A}$.

\section{Ultimate Logical Refinement}
True proofs of these structural problems can only be expressed using operations of ultimate logical refinement $\Omega$, defined such that for any finite composition $\phi$ of existing mathematical operations, $\phi \neq \Omega$. The limitation is thus not one of logical admissibility, but of absolute logical refinement:
\[
\forall \phi \in \mathrm{Finite} (\mathcal{A}), \phi \neq \Omega.
\]
Consequently, structural problems like the $\mathrm{ABC}$ conjecture lie beyond the operational scope of classical frameworks.

\section{Implications}
This formalization functions as both a conceptual and strategic guide: while existing frameworks present these problems as tractable conjectures, their intrinsic properties reveal that the true proofs reside outside the conventional operational vocabulary. This provides a mechanism to challenge assumptions regarding the solvability of structural problems using classical mathematics alone.

\section{Conclusion}
By framing classical mathematical frameworks as structural fixed points with minimal reasoning entropy, we clarify the inherent limitations in addressing complex structural problems. Structural conjectures such as the $\mathrm{ABC}$ conjecture demand operations of ultimate logical refinement, which exist outside the finite vocabulary of extant theories. This perspective can guide both the philosophical interpretation and the methodological approach to unresolved conjectures.

\end{document}
